% ID: SEL-D0046-OB-E04
% ESTADO: cotejada. Segunda pasada visual y borrador compilado revisados.
% Fuente: Selectividad/sel_2014_matematicas_aplicadas/Matematicas_CCSS_JUN.pdf
% SHA-256: df7375f8183eafb6f45e9c65083e68d528ffabe07d66717b0ccf0f020c5b064f
% Página: 2; opción B; ejercicio original 4.

\textbf{(2.5 puntos)} Un titular de prensa afirma que el 70\% de los jóvenes de una ciudad utilizan las redes sociales para comunicarse.

Para contrastar la veracidad de tal afirmación se toma una muestra aleatoria de 500 jóvenes de esa ciudad, y se obtiene que 340 de ellos utilizan la red para comunicarse.

Analice mediante un contraste de hipótesis bilateral, $(H_0:p=0.7)$, si se puede aceptar, con un nivel de significación del 1\%, que dicha afirmación es cierta.
